% CVPR Workshop 4-page format
% Target: Workshop on Generative Models for Creative Content, CVPR 2026
\documentclass[10pt,twocolumn,letterpaper]{article}

\usepackage{cvpr}
\usepackage{times}
\usepackage{epsfig}
\usepackage{graphicx}
\usepackage{amsmath}
\usepackage{amssymb}
\usepackage{booktabs}
\usepackage{multirow}
\usepackage{xcolor}
\usepackage[pagebackref=true,breaklinks=true,letterpaper=true,colorlinks,bookmarks=false]{hyperref}

\cvprfinalcopy  % comment out for blind review

\def\cvprPaperID{****}
\def\httilde{\mbox{\tt\raisebox{-.5ex}{\symbol{126}}}}

\begin{document}

\title{Captioning Strategy as a Determinant of Style Fidelity\\
in LoRA-based Manga Sub-genre Generation}

\author{An{\i}l Keskin\\
{\tt\small anil04keskin@gmail.com}
}

\maketitle

%------------------------------------------------------------------------
\begin{abstract}
Text-to-image diffusion models have demonstrated remarkable generative capability, yet they consistently reduce the rich diversity of manga into a homogeneous ``anime aesthetic.''
Fine-tuning via Low-Rank Adaptation (LoRA) offers a computationally tractable path toward sub-genre fidelity, but a critical upstream variable remains largely unexamined: the captioning strategy used to annotate training images.
We conduct a controlled study comparing two captioning approaches—BLIP-2 automatic captions (free-form natural language) and WD14 tagger outputs (structured Danbooru-style tags)—applied to a curated set of noir manga panels drawn from the Manga109-s dataset.
Both conditions are trained with identical LoRA hyperparameters on Stable Diffusion XL, and evaluated with FID, CLIP score, and DINOv2 style similarity against a held-out test set.
Our results show that BLIP-2 free-form captions produce stronger style fidelity (FID 261.97, DINOv2 similarity 0.322) while WD14 structured tags yield better text-prompt alignment (CLIP score 0.272 vs.\ 0.251), with implications for how practitioners should prioritize captioning strategies when fine-tuning on niche visual styles.
\end{abstract}

%------------------------------------------------------------------------
\section{Introduction}

The rapid maturation of latent diffusion models~\cite{rombach2022ldm,podell2023sdxl} has enabled photorealistic synthesis at scale, yet creative sub-genres remain stubbornly underserved.
Manga, and in particular its gekiga and noir variants, exhibits a coherent visual language—heavy ink linework, high-contrast screentone shading, claustrophobic panel compositions—that diverges sharply from the clean line art and pastel palettes associated with mainstream anime.
Prompt engineering alone fails to reproduce this language reliably, and subject-driven methods such as DreamBooth~\cite{ruiz2023dreambooth} or Textual Inversion~\cite{gal2022textual} are prohibitively expensive for researchers without access to large GPU clusters.

LoRA~\cite{hu2021lora} has emerged as the practical middle ground: it injects trainable low-rank matrices into frozen attention layers, achieving style adaptation in under an hour on a single A100.
However, the quality of a LoRA is fundamentally bounded by the quality of its training captions.
A caption that accurately names the visual elements of an image guides the cross-attention mechanism toward the correct text–image alignment; a generic or misleading caption can cause the LoRA to memorize spurious correlations rather than the intended style.
Despite the obvious importance of this choice, caption strategy has not been systematically isolated as an experimental variable in the fine-tuning literature.

This paper fills that gap.
We curate 200 noir manga panels from Manga109-s~\cite{matsui2017manga109}, apply two automated captioning strategies, train LoRA adapters under otherwise identical conditions, and evaluate style fidelity quantitatively and qualitatively.
Our contributions are:
\begin{itemize}
  \item A reproducible curation protocol for extracting style-coherent panels from annotated manga datasets.
  \item A direct, controlled comparison of BLIP-2~\cite{li2023blip2} and WD14 tag-based captioning as LoRA training supervision.
  \item An evaluation benchmark combining FID~\cite{heusel2017fid}, CLIP~\cite{radford2021clip}, and DINOv2~\cite{oquab2023dinov2} metrics for niche-style fidelity.
\end{itemize}

%------------------------------------------------------------------------
\section{Related Work}

\subsection{Fine-Tuning Diffusion Models for Style}

The dominant paradigm for personalizing large text-to-image models has evolved through three approaches.
DreamBooth~\cite{ruiz2023dreambooth} fine-tunes all model weights with a rare token bound to a target concept and employs a class-preservation loss to prevent language drift; this achieves strong subject fidelity but requires full model copies per style and is susceptible to overfitting on small datasets.
Textual Inversion~\cite{gal2022textual} instead fixes the model and optimizes a new token embedding, yielding a lightweight artifact but limited expressiveness because a single embedding vector must encode the full complexity of a novel concept.
LoRA~\cite{hu2021lora} offers a complementary trade-off: it injects trainable low-rank decompositions $\Delta W = BA$ (where $B \in \mathbb{R}^{d \times r}$, $A \in \mathbb{R}^{r \times k}$, $r \ll \min(d,k)$) into the query and value projections of each cross-attention layer, adding fewer than 4\% additional parameters while achieving style adaptation on par with full fine-tuning.
The expressiveness of LoRA scales with rank $r$; prior work has found rank 16–64 sufficient for style transfer tasks~\cite{hu2021lora}.
Critically, across all three paradigms, caption quality is treated as a given rather than a variable, motivating the present study.

\subsection{Captioning Strategies for Vision-Language Alignment}

BLIP-2~\cite{li2023blip2} bootstraps vision-language alignment by inserting a lightweight Querying Transformer between a frozen image encoder and a frozen large language model, generating fluent free-form captions with minimal fine-tuning cost.
While BLIP-2 produces grammatically natural descriptions, its training distribution—web-scraped image-text pairs—skews toward photorealistic content, making its captions less precise for niche artistic styles.
For a noir manga panel, BLIP-2 might output ``a black and white drawing of a man in a hat,'' capturing the denotative content but discarding the stylistic vocabulary (screentone, ink weight, compositional type) that the LoRA needs to internalize.

WD14 taggers~\cite{radford2021clip} take the opposite approach: they are trained on Danbooru, a large community-annotated illustration database, and output structured comma-separated tags drawn from a controlled vocabulary of several thousand terms.
These tags encode style-relevant attributes explicitly—\texttt{monochrome}, \texttt{screentone}, \texttt{dramatic lighting}, \texttt{1boy}, \texttt{hat}—and because Danbooru covers manga and anime at scale, the tagger's vocabulary is better matched to the domain.
The trade-off is that tags lack compositional structure; the cross-attention mechanism receives an unordered bag of terms rather than a syntactically coherent description, which may limit fine-grained alignment.

\subsection{Manga Generation and Style Transfer}

The Manga109 dataset~\cite{matsui2017manga109} provides bounding-box annotations for panels, text regions, character faces, and bodies across 109 manga volumes spanning decades of publication.
Its curated sub-split, Manga109-s, covers 87 works under an approved research license.
Prior work on manga generation has focused primarily on page layout synthesis and character consistency; DiffSensei~\cite{chen2024diffsensei} is among the first systems to address multi-panel manga story generation using a multimodal LLM to condition a diffusion backbone on character references and dialogue.
However, DiffSensei targets a broad range of styles and does not isolate sub-genre fidelity as a research objective.
To our knowledge, no prior work has examined how upstream captioning choices affect the style-specificity of LoRA-trained diffusion models on manga sub-genres.

%------------------------------------------------------------------------
\section{Method}

\subsection{Dataset Curation}

We extract panels from the 87 works in Manga109-s using the provided XML frame annotations, which give precise bounding boxes for each panel, text region, and character element.
Rather than relying on contour detection—which degrades on low-resolution scans—we use the ground-truth frame coordinates directly, extracting each panel at its native resolution and then resizing to $768 \times 768$ via center-crop with zero padding to preserve aspect ratio.

To select noir-style works, we inspect a stratified random sample of each title and retain works exhibiting at least two of the following markers: heavy use of black fills, screentone gradients, noirish subject matter (crime, urban, psychological thriller), and predominance of male adult characters in minimal-dialogue panels.
We sample up to 4 panels per work to ensure stylistic diversity without over-representing any single title.
After this process, we obtain 348 panels from 87 works, of which 200 are allocated to training and 40 to a held-out test set.
The train/test split is made at the \emph{work level}, not the page level, to prevent data leakage from a single artist's style appearing in both splits.

\subsection{Captioning Conditions}

All captions are prefixed with the trigger token \texttt{gknoir style,} to bind the LoRA to a unique, out-of-vocabulary identifier.

\paragraph{Condition A — BLIP-2.}
We load \texttt{Salesforce/blip2-opt-2.7b} in 8-bit precision and run inference on each training panel with greedy decoding.
The raw caption is prepended with the trigger, yielding descriptions such as ``\textit{gknoir style, a black and white drawing of a man in a coat walking down a rainy street.}''

\paragraph{Condition B — WD14.}
We run \texttt{SmilingWolf/wd-v1-4-convnext-tagger-v2} via its ONNX export, applying a confidence threshold of 0.35 and filtering copyright tags (Danbooru category 9).
The output tag set is sorted by confidence and concatenated, then prepended with the trigger: ``\textit{gknoir style, 1boy, monochrome, hat, dramatic lighting, screentone, noir.}''

\subsection{LoRA Training}

We train both conditions on \texttt{stabilityai/stable-diffusion-xl-base-1.0} using kohya\_ss sd-scripts with identical hyperparameters: rank $r=32$, $\alpha=16$, AdamW8bit optimizer, learning rate $10^{-4}$ with cosine-with-restarts schedule (100 warmup steps), 2000 training steps, batch size 2, \texttt{bf16} mixed precision, and xformers attention.
Each work's panels are repeated $10\times$ per epoch via the \texttt{N\_classname} directory convention.
Training runs on a single A100 40\,GB GPU and completes in approximately 45 minutes per condition.

\subsection{Evaluation Protocol}

We generate 100 images per condition using 10 fixed prompts (Table~\ref{tab:prompts}) crossed with 10 seeds, and compute three metrics.

\textbf{FID}~\cite{heusel2017fid} measures the Fr\'{e}chet distance between Inception-v3 feature distributions of the 100 generated images and the 40 held-out test panels.
Lower FID indicates that the generated distribution is closer to the real manga distribution.

\textbf{CLIP score}~\cite{radford2021clip} computes the cosine similarity between each generated image's CLIP-ViT-L/14 embedding and its conditioning prompt embedding, averaged across all 100 images.
Higher scores indicate better text–image alignment.

\textbf{DINOv2 style similarity}~\cite{oquab2023dinov2} extracts \texttt{dinov2-vitl14} patch-level features from generated images and test panels, and computes the mean cosine similarity between the two distributions.
DINOv2 features are more sensitive to low-level texture and style than Inception or CLIP features, making this metric especially informative for stylistic fidelity assessment.

We also report all three metrics for a \textbf{baseline} condition: 100 images generated from the same prompts using SDXL with no LoRA.

%------------------------------------------------------------------------
\section{Experiments}

\subsection{Quantitative Results}

\begin{table}[h]
\centering
\caption{Quantitative evaluation results. FID lower is better; CLIP score and DINOv2 similarity higher is better.}
\label{tab:quant}
\begin{tabular}{lccc}
\toprule
\textbf{Condition} & \textbf{FID}$\downarrow$ & \textbf{CLIP}$\uparrow$ & \textbf{DINO}$\uparrow$ \\
\midrule
Baseline (no LoRA) & 306.46 & \textbf{0.2773} & 0.155 \\
BLIP-2 (Cond.\ A)  & \textbf{261.97} & 0.2508 & \textbf{0.3219} \\
WD14   (Cond.\ B)  & 272.13 & 0.2718 & 0.2738 \\
\bottomrule
\end{tabular}
\end{table}

\subsection{Qualitative Comparison}

\begin{figure}[h]
  \centering
  % \includegraphics[width=\linewidth]{figures/qual_grid.png}
  \fbox{\parbox{0.9\linewidth}{\centering \vspace{2cm} \textit{Qualitative comparison grid placeholder.\\
  Rows: 3 eval prompts. Cols: Baseline / BLIP-2 / WD14.} \vspace{2cm}}}
  \caption{Side-by-side generation comparison for a fixed seed across all three conditions.}
  \label{fig:qual}
\end{figure}

\subsection{Evaluation Prompts}

\begin{table}[h]
\centering
\caption{Fixed evaluation prompts used for all conditions.}
\label{tab:prompts}
\resizebox{\linewidth}{!}{%
\begin{tabular}{cl}
\toprule
\textbf{\#} & \textbf{Prompt} \\
\midrule
1 & gknoir style, heavy ink linework, close-up, detective in rain, tense \\
2 & gknoir style, thin linework, wide shot, empty city street, melancholic \\
3 & gknoir style, heavy ink, over-shoulder, figure in doorway, stark shadows \\
4 & gknoir style, bold ink, action panel, two figures confronting, tension \\
5 & gknoir style, fine linework, interior, desk with lamp, quiet solitude \\
6 & gknoir style, heavy ink, extreme close-up, eyes in shadow, mysterious \\
7 & gknoir style, crosshatch shading, mid-shot, woman in coat, rainy street \\
8 & gknoir style, stark contrast, running figure, urgency \\
9 & gknoir style, screentone, thoughtful expression, window light, introspective \\
10 & gknoir style, deep blacks, alley at night, architecture detail, foreboding \\
\bottomrule
\end{tabular}}
\end{table}

%------------------------------------------------------------------------
\section{Discussion and Conclusion}

\paragraph{Findings.}
Table~\ref{tab:quant} reveals a consistent but nuanced pattern.
Both LoRA conditions substantially reduce FID relative to the baseline (BLIP-2: $-$44.5, WD14: $-$34.3), confirming that fine-tuning on even a small curated set markedly shifts the generated distribution toward the target style.
BLIP-2 achieves the lower FID (261.97 vs.\ 272.13) and the higher DINOv2 similarity (0.3219 vs.\ 0.2738), indicating stronger \emph{style} fidelity under both distributional and feature-level measures.
WD14 recovers on text-prompt alignment: its CLIP score of 0.2718 is close to the baseline (0.2773), while BLIP-2 drops to 0.2508.
The answer to our research question is therefore: BLIP-2 free-form captions produce a LoRA that is more faithful to the source style; WD14 tags produce a LoRA that follows structured evaluation prompts more precisely.

\paragraph{Why BLIP-2 wins on style.}
This result initially appears counterintuitive.
WD14 tags supply domain-specific vocabulary (\texttt{screentone}, \texttt{monochrome}, \texttt{dramatic lighting}) absent from BLIP-2's photorealistic training distribution, so one might expect WD14 to better communicate the target aesthetic.
However, the key variable is \emph{coverage}, not domain specificity.
BLIP-2 captions describe the full visual composition of each panel — spatial relationships, mood, character posture — providing richer cross-attention targets for every image.
WD14 tags, being an unordered bag of terms, offer no compositional structure; the model receives the same set of words regardless of layout.
We hypothesize that BLIP-2's holistic descriptions allow the LoRA to internalize not just the isolated visual atoms of noir manga (ink weight, screentone) but their characteristic arrangement, which is precisely what DINOv2 — sensitive to spatial feature distributions — rewards.
The CLIP score trade-off is consistent with this explanation: WD14 tags overlap more directly with the structured evaluation prompts (also comma-separated stylistic terms), inflating cosine similarity without necessarily improving style reproduction.

\paragraph{Limitations.}
Our dataset (200 training panels from 87 works) is small relative to typical LoRA corpora; conclusions may not hold at larger scale.
We study a single sub-genre of manga, and it is unclear whether BLIP-2's compositional advantage generalizes to genres with more uniform panel layouts.
Our FID estimate uses 100 generated images against 40 test images, below the 10,000-image sample recommended for reliable FID~\cite{heusel2017fid}; all FID values should therefore be interpreted as indicative rather than definitive.

\paragraph{Conclusion.}
We have presented a controlled study isolating captioning strategy as a determinant of LoRA style fidelity on noir manga.
BLIP-2 free-form captions produce higher style fidelity (FID, DINOv2) while WD14 tags yield better text-prompt alignment (CLIP), and practitioners should select the captioning strategy based on whether style reproduction or prompt responsiveness is the primary objective.
We release all curation code, caption files, training configs, and evaluation scripts to support reproducibility and future comparisons.

%------------------------------------------------------------------------
{\small
\bibliographystyle{ieee_fullname}
\bibliography{references}
}

\end{document}
